\section{Yoneda structure}
\label{sec:yoneda}

We'll need universal properties in the $2$-category of
multi-categories. It might be useful to develop a rich theory of
representability inside $\MultiCAT$. We'll try to do it using the
\emph{Yoneda structure} of \cite{street-walters-yoneda-structures}.

\subsection{Admissible multi-functors}
A multi-functor $F : \UU \to \UU[2]$ is \emph{admissible} when, for
all $a_1, \ldots, a_n$ in $\UU$ and $b$ in $\UU[2]$, the multi-homset
$\UU[2](Fa_1, \ldots, Fa_n; b) \in \Set$ is small. A multi-category is
\emph{admissible} when its identity multi-functor is admissible, i.e.,
when it is locally small. The admissible multi-functors are closed
under-precomposition, i.e., for $\UU[1] \xto F \UU[2] \xto G \UU[3]$,
if $G$ is admissible, then $G \compose F$ is also admissible. Indeed,
given $a_i$ in $\UU[1]$ and $c$ in $\UU[3]$, we have:
\[
\UU[3](G\compose F a_1, \ldots, G\compose F a_n; C)
=
\UU[3](G(F a_1), \ldots, G(F a_n); C)
\explain \in {$G$ admissibility\\for $Fa_1, \ldots, Fa_n$}
\Set
\]

\subsection{Iterated lists}
As you can see in from previous sections, manipulating
multi-categories in the abstract makes the indexing heavy. Here we
develop some notation to manipulate sequences of multi-morphisms
ensemble, which is probably related to a construction of a free
monoidal category out of a multi-category.

Recall the structure of the list/free monoid monad, given by:
\begin{itemize}
  \item For each set $X$, the set $\List X$ of $X$-sequences $\xvec :=
    \seq{\xvec_1, \ldots, \xvec_{\abs\xvec}}$, i.e., $\abs\xvec$ is
    the length of the sequence, and its components in order are
    $\xvec_i$ for $1 \leq i \leq n$.

    We will also grade the monad, i.e., $\List_n X \subset \List X$
    are the length $n$ sequences.

  \item For each $x \in X$, the singleton $\seq{x} \in \List_1 X$.
  \item Given a sequence $\xvec \in \List_n X$ and, for each $1 \leq i
    \leq \abs\xvec$, another sequence $\yvec_i \in \List_{m_i}Y$, we
    have the concatenated sequence:
    \[
    \concat_{i=1}^{n}\yvec^i := \seq{\yvec_{11}, \ldots,
      \yvec_{1m_1}, \ldots, \yvec_{n1}, \ldots,
      \yvec_{nm_n}}
    \]
    We will also need iterated lists. These too form a monad-like
    structure, that of \emph{unlabelled, rooted, finitely branching
      trees of uniform depth}:
    \begin{itemize}
    \item For each set $X$, the set $\Tree X \definedby \Union_{k \in
      \naturals} \List^kX$ represents finitely-branching rooted trees
      whose nodes are unlabelled, and whose leaves, all of which have
      the same depth, are taken from $X$.
    \item For each $x \in X$, we identify the unique rooted tree with
      $1$ leaf at depth $1$ labelled by $x$ with the singleton list
      $\seq x \in \List^1X$.

      The empty tree, that only has the root node and no leaves, is
      represented by the empty list.

      Note that trees are allowed to have childless \emph{nodes} at
      non-uniform depths, e.g.: $\seq{\seq{},\seq{\seq{}},
        \seq{\seq{\seq{1}}}}$ has two nodes at depths $1$ and $2$ and
      a leaf at depth $3$.
    \item Given a tree whose leaves are trees, we can substitute the
      inner trees into the outer tree.
    \end{itemize}

    Let $[n] := \set{1, \ldots, n}$ be the canonical set of $n$-many
    indices.  For each depth-$k$ iterated list $\xivec \in \List^k X$,
    let $\floor\xivec \in \List X$ be its fully flattened sequence.

    We denote the set of \emph{branches} of a rooted tree $\xivec$ by
    $\Branch \xivec$, i.e., sequences of indices in the tree that mark
    a path from the root to a leaf. The leaf-depths are uniform, we
    have, for each $\xivec \in \List^kX$:
    \[
    \Branch \xivec = \coprod_{i_1 \in
      [\abs\xivec]}\coprod_{i_2 \in [\abs{\xivec_{i_1}}]}\cdots
    \coprod_{i_{k-1}\in[\abs{\xivec_{i_1,i_2\ldots,i_{k-2}}}]}[\abs{\xivec{i_1,
          \ldots, i_{k-1}}}]
    \]

    Let the \emph{fringe size} of a tree be its number of leaves,
    i.e., the cardinality of its set of branches:
    \begin{align*}
    \norm \xivec \definedby \abs{\floor\xivec}
    &{}\definedby \sum_{i_1 \in [\abs\xivec]} \sum_{i_2 \in [\abs{\xivec_{i_1}}]}\cdots
    \sum_{i_{k-1} \in [\abs{\xivec_{i_1,i_2,\ldots,i_{k-2}}}]}\abs{\xivec_{i_1,i_2,\ldots,i_{k-2},i_{k-1}}}
    \\&{}=\cardinality{\Branch \xivec}
    \end{align*}
    For example, for $\xivec \definedby
    \seq{\seq{\seq{1,2},\seq{3,4}},\seq{\seq{5},\seq{6,7},\seq{8}}}$
    we have $\floor\xivec = \seq{1,2,\ldots,8}$ and:
    \[
    \abs{\xivec} = 2\qquad
    \abs{\xivec_1} = 2 \qquad
    \abs{\xivec_2} = 3 \qquad
    \norm{\xivec} :=
    \sum_{i_1 \in [\abs\xivec]}\sum_{i_2 \in [\abs{\xivec_{i_1}}]}\abs{\xivec_{i_1,i_2}}
    = \parent{2 + 2} + \parent{1 + 2 + 1} = 8
    \]
    We then define the bijection:
    \begin{align*}
    \seq[{\xivec}]{-,-,\ldots,-} &{}: \Branch\xivec
    \xto{\isomorphic}
        [\norm\xivec]
    \\\seq[\xivec]{i_1, \ldots, i_k} &{}\definedby
    \sum_{j_1 < i_1}\norm{\xivec_{j_1}}
    +
    \sum_{j_2 < i_2}\norm{\xivec_{i_1,j_2}}
    +
    \ldots
    +
    \sum_{j_{k-1} < i_{k-1}}\norm{\xivec_{i_1,i_1, \ldots, i_{k-2},j_{k-1}}}
    + i_k
    \end{align*}
    so that $\floor\xivec_{\seq{i_1, \ldots, i_k}} = \xivec_{i_1,
      \ldots, i_k}$ and $\seq{\xivec_{\seq[\xivec][-1]{1}}, \ldots,
      \xivec_{\seq[\xivec][-1]{\norm\xivec}}}$. For example,
    continuing with the previous $\xivec$, we have:
    \begin{align*}
      \seq{1,2,1} &= 0 + 2 + 1 = 3 & \floor\xivec_3 &= 3 = \xivec_{1,2,1}
      \\
      \seq{2,2,1} &= 4 + (1 + 2) + 1 = 8 & \floor\xivec_8 &= 8 = \xivec_{2,2,1}
    \end{align*}
\end{itemize}
Given a list $\avec \in \List A$ over some set $A$, an \emph{ordered
  partition of $\avec$} is a list-of-lists $\pi \in \List^2 A$ such
that $\concat_{i = 1}^{\abs \pi} \pi_i = \avec$. Similarly, an
\emph{iterated ordered partition of $\avec$} is a tree $\pi \in \Tree
A$ such that $\floor \pi = \avec$.


\subsection{The multi-category of presheaves}
Let $\UU$ be a multi-category. We will define a \emph{category}
$\opposite\UU$ so that our presheaves will be functors $F :
\opposite\UU \to \Set$:
\begin{itemize}
\item $\Obj{\opposite\UU} := \List \Obj\UU$ are finite-sequences
  $\avec = \seq{\avec_1, \ldots, \avec_n}$ of $\UU$-objects.
\item Morphisms $f \in \opposite\UU(\avec , \bvec)$, which we will
  write as ``$f : \avec \from \bvec$ in $\UU$'', consist of:
  \begin{itemize}
  \item an ordered partition $\partit_f\definedby \seq{\partit_1,
    \ldots, \partit_n} \in \List^2 \Obj\UU$, $\vec b$, i.e.:
    \(
    \bvec = \concat_{i=1}^n \partit_i
    \).
  \item for each $1 \leq i \leq n$, a multi-morphism $f_i : \avec_i \from
    (\bvec_i)$.  We can depict a morphism like so:
    \[
    \quiver{Cop morphism}
    \]
  \item The identity $\id[\vec a]^{[\opposite\UU]}$ consists of:
    \begin{itemize}
    \item $\partit_{\id[\vec a]} := \seq{\seq{a_1}, \ldots,
      \seq{a_n}}$, the partition into singletons; and
    \item $(\id[\vec a])_i := \id[a_i] : a_i \from (a_i)$.
    \end{itemize}
  \item To compose morphisms $\avec \xfrom f \bvec \xfrom g \cvec$,
    use the multiplication of the $\Tree$ monad to get an iterated
    partition of depth $3$, and then flatten all the depth $1$
    iterated lists into lists.  Then, compose the multi-morphisms to
    get a morphism $f \compose g : \avec \from \cvec$ (note the order
    of arguments in this notation). Letting $\abs\avec \defines n$ and
    $\abs{\partit_{f,i}} \defines m_i$:
    \begin{itemize}
    \item $\partit_{f \compose g} \definedby
      \seq{\seq{\partit_{g_{\pair11}}
                \++ \ldots \++
                \partit_{g_{\pair1{m_1}}}}
          , \ldots
          ,\seq{\partit_{g_{\pair n1}}
                \++ \ldots \++
                \partit_{g_{\pair n{m_n}}}}
          }$
    \item $(f \compose g)_{i} := f_i\compose (g_{\pair i1}, \ldots, g_{\pair i{m_i}})$.
    \end{itemize}
      We can visualise this operation so, letting $k_{i,j} \defines
      \abs{\partit_{g_i},j}$:
      \[
      \quiver{Cop composition}
      \]
  \end{itemize}
\end{itemize}
\begin{lemma}
  This data defines a category $\opposite\UU$.
\end{lemma}
\begin{proof}
  We inherit the properties from the monad laws and the multi-category
  axioms.
\end{proof}

A \emph{presheaf over $\UU$} is a functor $F : \opposite\UU \to \Set$.

\begin{example}\examplelabel{Yoneda mapping}
  Let $\UU$ be a locally small multi-category. For each $a \in
  \Obj\UU$, we have a presheaf $\yoneda[\UU]a : \opposite\UU \to \Set$
  called the presheaf \emph{represented by $a$}:
  \[
  \yoneda_{\UU}\,a\,\bvec \definedby \UU(\bvec; a) \in \Set
  \qquad
  \yoneda_{\UU}\,a\,(\bvec \xfrom{f} \cvec) : g \in \UU(\bvec; a)
  \mapsto g \compose (f_1, \ldots, f_{\abs\bvec}) \in \UU(\cvec; a)
  \]
  The functoriality:
  \[
  \yoneda_{\UU}\,a\,(\bvec \xfrom{f} \cvec \xfrom{g} \dvec)
  =
  \yoneda_{\UU}\,a\,(g)
  \compose
  \yoneda_{\UU}\,a\,(f)
  \qquad
  \yoneda_{\UU}\,a\,(\id[\bvec])
  =
  \id[\yoneda_{\UU}\,a\,\bvec]
  \]
  follows from the multi-category axioms and the definition of $\opposite\UU$.
\end{example}

Let $F_1, \ldots, F_n, G : \opposite\UU \to \Set$ be presheaves.  A
\emph{natural multi-transformation $\alpha : (F_1, \ldots, F_n) \to
  G$} consists of:
\begin{itemize}
\item For every sequence $\avec\in \Obj{\opposite\UU}$ and length-$n$
  ordered partition $\partit$ of $\avec$, a function:
  \[
  \alpha_{\partit} : \prod_{i=1}^nF_i\partit_i \to G\avec
  \]
\end{itemize}
such that:
\begin{itemize}
\item For every $f : \avec \from \bvec$ in $\UU$ and for every
  length-$n$ ordered partition $\phi = \seq[i=1][n]{\phi_i : \partit_i
    \from \partit'_i}$ of the sequence $f$: \diagram{19}
\end{itemize}

\begin{example}\examplelabel{Yoneda functoriality}
  Let $g : (\avec) \to b$ be a multi-morphism in $\UU$. We define a
  natural multi-transformation, where $n \definedby \abs\avec$:
  \begin{align*}
  \yoneda_{\UU}\,g :{}& (\yoneda_{\UU}\, \avec_1, \ldots, \yoneda_{\UU}\,\avec_n) \to
  \yoneda_{\UU}\,b\\
  (\yoneda_{\UU}\,g)_{\partit} :{}&
  \seq{ \partit_1 \xto{u_1} \avec_1
      , \ldots
      , \partit_n \xto{u_n} \avec_n
      } \mapsto
  g\compose (u_1, \ldots, u_n)
  \end{align*}
  To see that it is natural, take any $f : \cvec \from \dvec$ in
  $\UU$, and an $n$-ary ordered partition of it $\phi$, and chase
  elements around the naturality diagram:
  \diagram{20}
  and so $\yoneda g$ is a natural multi-transformation.
\end{example}
The multi-category $\Psh\UU$ of \emph{presheaves over $\UU$} consists
of this data:
\begin{itemize}
\item Objects are presheaves $F : \opposite\UU \to \Set$.
\item Multi-morphisms $\alpha : (F_1, \ldots, F_n) \to G$ are the
  natural multi-transformations.
\item The identity multi-transformation $\id[F] : (F) \to F$ has as
  component at each unary partition $\partit = \seq{\avec}$ of $\avec$
  the canonical isomorphism:
  \[
  (\id[F])_{\seq{\avec}} : \prod\set{F\avec} \xto{\isomorphic} F\avec
  \]
\item To define composition, take composable natural
  multi-transformations
  \[
  \alpha = \seq[i=1][n]{(\Fvec_i) \to \Gvec_i}
  \quad\text{and}\quad
  \beta : (\Gvec) \to \Hvec
  \quad\text{where:}\quad
  \Fvec \in \List^2(\opposite\UU \to \Set),
  \Gvec \in \List^1(\opposite\UU \to \Set)
  \] letting as usual $p := \norm
  \Fvec$, $n \definedby \abs \Gvec$, and $m_i \definedby
  \abs{\Fvec_i}$. Given a $p$-ary ordered partition $\partit$ of some
  $\avec$, define these two auxiliary partitions:
  \begin{itemize}
  \item $\partit_{\pair i*} \definedby \seq[j=1][m_i]{\partit_{\pair ij}}$
  \item $\partit_{\pair i{\++}} \definedby
    \concat_{j=1}^{m_i}\partit_{\pair ij} = \floor{\partit'_i}$
  \item So $\alpha_{i, \partit_{\pair i*}}
    :\prod_{j=1}^{m_i}\Fvec_{\pair ij} \to \Gvec_i\partit_{\pair
    i{\++}}$.
  \end{itemize}
  It might easier to visualise how they relate to each other by
  aligning their fringes:
  \[
  \begin{array}{@{}*{4}{l@{}}c@{}l@{}c@{}l@{}@{}l@{}c@{}l@{}*{2}{l@{}}}
    \partit {}&= \tabmark
                  \iseq{\tabmark         \partit_{\pair 11}\
                       \tabmark,\ldots, \tabmark\partit_{\pair 1{m_1}}
                       \tabmark,\ldots,\tabmark\tabmark\partit_{\pair n1}
                       \tabmark,\ldots, \tabmark\partit_{\pair n{m_n}}\tabmark}
    \\
    \partit_{\pair{*}*} &{}:=\tabmark
                  \iseq{\iseq{\tabmark  \partit_{\pair 11}
                             \tabmark,\ldots, \tabmark\partit_{\pair 1{m_1}}}
                             \tabmark,\ldots,\tabmark
                        \iseq{\tabmark\partit_{\pair n1}
                             \tabmark,\ldots,\tabmark\partit_{\pair n{m_n}}}\tabmark}
    \\
    \partit_{\pair{*}{\++}} &{}:=\tabmark
                  \iseq{\tabmark               \partit_{\pair 11}
                        \tabmark\++ \ldots\++\tabmark\partit_{\pair 1{m_1}}
                        \tabmark,\ldots,\tabmark\tabmark\partit_{\pair n1}
                        \tabmark\++ \ldots\++\tabmark\partit_{\pair n{m_n}}
                  \tabmark}
  \end{array}
  \]
  We define the $\partit$-th component of the composition:
  \[
  (\beta \compose (\alpha_1, \ldots, \alpha_n))_{\partit} :
  \prod_{\ell =1}^p\Fvec_{\ell}\partit_{\ell} \xto{\isomorphic}
  \prod_{i = 1}^n\prod_{j=1}^{m_i} \Fvec_{\pair ij}\partit_{\pair ij}
  \xto{\prod_{i=1}^n\alpha_{i,\partit_{\pair i*}}}
  \prod_{i=1}^n\Gvec_i\partit_{\pair i{\++}}
  \xto{\beta_{\partit_{\pair *{\++}}}}
  H\avec
  \]
\end{itemize}
\begin{proposition}
  These data define a multi-category $\Psh\UU$.
\end{proposition}
\begin{proof}
  First, we need to show the identities and composition are
  well-defined, i.e., natural. For the identity, take any unary
  partition $\phi = \seq{f}$ of a morphism $f : \avec \from \bvec$ in
  $\UU$, and calculate:
  \diagram{21}
  Turning to the well-definedness of composition, continue with the
  same notation and take a $p$-ary partition $\phi = \seq{\phi_{\ell} :
    \partit[1]_{\ell} \from \partit[2]_{\ell}}$ of some morphism $f : \avec \from
  \bvec$ in $\UU$. The crucial step is to repeat the same partitioning
  scheme for $\phi$ as well:
  \begin{itemize}
  \item $\phi_{\pair i*} \definedby \seq[j=1][m_i]{\phi_{\pair ij} : \partit_{\pair ij} \from \partit[2]_{\pair ij}}$;
  \item $\phi_{\pair i{\++}} \definedby \concat_{j=1}^{m_i}\phi_{\pair ij} :
    \partit_{\pair i{\++}} \from \partit[2]_{\pair i{\++}}$ so that $\phi_{\pair i*}$ is a partition of $\phi_{\pair i{\++}}$;
  \end{itemize}
  and we can now calculate the naturality condition of the composition:
  \diagram{22}

  The neutrality of the identities w.r.t.~composition follow
  immediately:
  \begin{center}
  \diagram*{23}
  \diagram*{24}
  \end{center}

  To prove that composition is associative, we'll first dwell on the
  partitioning scheme we used to define composition.

  We can organise the data in a sequence of $d$-many composable
  multi-transformations:
  \[
  \Fvec^{d+1} \xto{\alpha^{d}} \cdots \xto{\alpha^1} \Fvec^1 \xto{\alpha^0} F^0
  \tag{($*$)}
  \]
  as follows. We consider a rooted tree $t$ of uniform leaf-depth $d+1$. We label
  each leaf, say in branch $\seq{i_1, \ldots, i_{d+1}}$ in the tree,
  with a single presheaf $\Fvec^{d+1}_{\seq{i1, \ldots,
      i_{d+1}}}$. Each node of depth $1 \leq s \leq d$ in the tree,
  say on the path $\seq{i_1, \ldots, i_s}$ is also labelled with a
  presheaf $\Fvec^s_{i_1, \ldots, i_s}$. Given these labels, we also
  annotate each node, that has, say, $k_{\seq{i_1, \ldots, i_s}}$
  children, with a natural multi-transformation $\alpha^s_{\seq{i_1,
      \ldots, i_s}} : \parent{\seq[i=1][k_{\seq{i_1, \ldots,
          i_s}}]{\Fvec^{s+1}_{i_1, \ldots, i_{s + 1}}}} \to
  \Fvec^s_{\seq{i_1, \ldots, i_s}}$.

  The bijection between the set of length-$s$ branches and the set of
  nodes and leafs at level $s$ allows us to define the lists
  $\Fvec^s$ and $\alpha^s$ to make the notation in ($*$) precise.
  After composing these multi-transformations in some order, we get a
  multi-transformation:
  \[
  \bigcirc \alpha: \parent{\Fvec^{s+1}} \to \Fvec^0
  \]
  Letting $p \definedby \abs{\Fvec^{s+1}}$, when we take any $\avec$
  and a $p$-ary partition of $\avec$, we can extend the tree $t$ by
  another level (turning all the previous leaves into nodes but
  without the multi-transformation label), and the leaves in the newly
  added level are labelled by the objects in $\avec$. It is this tree,
  its labels, and its branches and paths that we are manipulating with
  our notation.

  Given a tree (with any sort of labelling) of uniform leaf depth $d$,
  we denote by $t_{\seq{i_1, \ldots, i_s, *, \ldots, *}}$ the sub-tree
  rooted at node $\seq{i_1, \ldots, i_s}$. It's perhaps convenient to
  pad the remaining components up to a length-$d$ sequence to spot
  typos. We denote by:
  \[
  t_{\seq*{i_1, \ldots, i_{s},
      *, \ldots, *, \explain*{\++}{pos.~$j_1$},
      *, \ldots, *, \explain*{\++}{pos.~$j_2$},
     *,
     \underset{\substack{~\\[8pt]\ldots}}\ldots, *, \explain*{\++}{pos.~$j_r$},
     *, \ldots, *
  }}
  \]
  the sub-tree rooted at $\seq{i_1, \ldots, i_s}$, whose levels $j_1,
  \ldots, j_r$ were flattened. This notation is consistent with the
  notation we used for composition and in the previous parts of the
  proof, apart from identifying the partition $\pi$ with its depth-$3$
  iterated list. The crux of our proof is a kind of substitution
  lemma/monad-like associativity law that means the various iterations
  on this notation are consistent with each other, for example:
  \[
  (t_{i_1, i_2, *, \++, *, \++, *})_{j_1, \pair{j_2}{j_3}, \++}
  =
  t_{i_1, i_2, j_1, j_2, j_3, \++, \++}
  \]

  Back to proving associativity. Consider $3$ levels of composable
  multi-transformations:
  \[
  \Fvec \xto{\alpha} \Gvec \xto{\beta} \Hvec \xto{\gamma} \Kvec
  \]
  i.e.:
  \begin{itemize}
  \item $\gamma : \parent{\seq[i=1][k_{\seq{}}]{\Hvec_{\seq i}}} \to \Kvec$
  \item $\beta_{\seq i} :
    \parent{\seq[j=1][k_{\seq{i}}]{\Gvec_{\seq{i,j}}}} \to \Hvec_{\seq
      i}$, for all $1 \leq i \leq k_{\seq{}}$.
  \item $\alpha_{\seq{i,j}} :
    \parent{\seq[r=1][k_{\seq{i,j}}]{\Fvec_{\seq{i,j,r}}}} \to
    \Gvec_{\seq{i,j}}$, for all $1 \leq i \leq k_{\seq{}}$ and $1 \leq
    j \leq k_{\seq{i}}$.
  \end{itemize}
  Let:
  \begin{itemize}
  \item $f := \abs\Fvec$
  \item $f_{\seq{i,j}} := \abs{\Fvec_{\triple ij*}}$
  \item $g := \abs\Gvec$
  \item $g_i := \norm{\Gvec_{\seq{i, *}}} = k_{\seq i}$
  \item $h := \abs\Hvec = k_{\seq{}}$
  \end{itemize}
  Consider any $f$-ary partition $\partit$ of some $\avec$, and
  calculate:
  \diagram{25}
  Therefore, composition is associative.
\end{proof}

\begin{lemma}
  When $\UU$ is small, $\Psh\UU$ is locally small.
\end{lemma}
\begin{proof}
  When $\Obj\UU$ is small, so is $\List\Obj\UU$, and so is the set
  $\coprod_{\avec \in \List\Obj\UU}\Partit_n\avec$ of $n$-ary ordered
  partitions over some list of objects. Given a presheaf sequences and
  a presheaf $\Fvec,G : \UU \to \Set$, the collection of natural
  multi-transformation is a subset of the set $\prod_{\avec \in
    \Obj{\opposite\UU}, \partit\in\Partit\avec}
  \Set(\prod_i\Fvec_i\partit_i, G\avec)$.
\end{proof}

\subsection{The Yoneda lemma}
Even if it's not immediately clear from the axiomatisation, there's a
Yoneda lemma sitting at the heart of the Yoneda structure axioms.
Recall from examples \exampleref{Yoneda mapping} and
\exampleref{Yoneda functoriality} that, for every locally small
multi-category $\UU$, we have a multi-functor structure:
\[
\begin{array}{@{}l@{}l@{}}
  \yoneda_{\UU} &{}: \UU \to \Psh\UU \\
  \yoneda_{\UU}\,a &{}\definedby \UU(-,a) \\
  (\yoneda_{\UU}\,\parent{\vto \avec f{a'}})_{\partit} &{}\definedby
    \seq[i=1][\abs\avec]{(\partit_i)
    \xto{u_i} \avec_i} \mapsto f \compose \seq[i=1][\abs\avec]{u_i}
\end{array}
\]

\begin{lemma}
  This data defines a multi-functor $\yoneda_{\UU} : \UU \to \Psh\UU$.
\end{lemma}
\begin{proof}
  For identity multimorphisms $\id[a] : (a) \to a$, there is only one
  $1$-partition of any sequence $\bvec$, and we have:
  \[
  (\yoneda\id)_{\seq[i=1][1]{\bvec}} \seq[i=1][1]{u : (\bvec) \to a}
  \mapsto \id \compose (u) = u = (\id[\yoneda_a])_{\seq\bvec}\seq u
  \]
  and so $\yoneda\id = \id[\yoneda_a]$. Taking composable
  multi-morphisms $g : (\cvec) \to d$ and $f_i : (\bvec_i) \to
  \cvec_i$, $i = 1, \ldots, n := \abs\cvec$, take any $p$-ary
  partition $\partit$ of some $\avec$, where $p := \norm\bvec$, and
  chase an arbitrary element of $\prod_{\ell=1}^p\yoneda
  \bvec_\ell\partit_{\ell}$, where $m_i := \abs{\bvec_i}$:
  \diagram{26} and so $\yoneda(g \compose \seq[i]{f_i}) = \yoneda g
  \compose \seq[i]{\yoneda f_i}$.
\end{proof}

\begin{lemma}\emph{Yoneda for multi-categories.}
  Let $\UU$ be a locally-small multi-category, $\avec$ an $n$-ary
  sequence of objects in $\UU$, and $F \in \Psh\UU$ a presheaf. We
  have a bijection:
  \[
  \Yoneda_{\UU, F, \avec} : \Psh\UU(\yoneda \avec_1, \ldots, \yoneda\avec_n; F)
  \xto\isomorphic F\avec
  \]
  given by:
  \begin{gather*}
  \Yoneda_{\UU, F, \avec}(\seq[i=1][n]{\yoneda \avec_i} \xto\alpha F)
  \definedby
  \alpha_{\seq[i=1][n]{\seq[j=1][1]{\avec_i}}}\seq[i=1][n]{\id[a_i]}
  \in F\avec
  \\
  (\inv\Yoneda_{\UU, F, \avec}x)_{\partit} : \seq[i=1][n]{\partit_i
    \xto{u_i} \avec_i} \mapsto F(u_1, \ldots, u_n)x
  \end{gather*}
\end{lemma}
\goodbreak
\begin{proof}
  As usual, the proof is to just do it.

  To show well-definedness, i.e., that $\inv\Yoneda x : (\yoneda
  \avec_1, \ldots, \yoneda\avec_n) \to F$ is natural, take any
  morphism $\bvec \xfrom f\cvec$ in $\UU$, an $n$-ary partition
  $\seq[i=1][n]{\partit_i \xfrom{\phi_i} \partit'_i}$ of it, let
  $m_i:= \abs{\partit_i}$, and chase an arbitrary element of
  $\prod_{i=1}^n\yoneda \avec_i \partit_i$:
  \diagram{27}
  To show $\Yoneda \compose \inv\Yoneda = \id$, take any $x \in
  F\avec$:
  \[
  \Yoneda(\inv\Yoneda x) = F(\id, \ldots, \id)(x) \explain={functoriality} x
  \]
  For the converse, take any natural multi-transformation $\alpha :
  (\yoneda\avec_1, \ldots, \avec_n) \to F$. We need to show that, for
  every $n$-ary partition $\partit$ of some $\bvec$ and $\partit_i
  \xto{u_i} \avec_i$:
  \[
  \alpha_{\partit}(u_1, \ldots, u_n)
  =
  (\inv\Yoneda (\Yoneda \alpha))_{\partit}(u_1, \ldots, u_n) =
  F(u_1, \ldots, u_n)
  \parent{\alpha_{\seq[i=1][n]{\seq[j=1][1]{\avec_i}}}
    (\id[\avec_1], \ldots, \id[\avec_n])}
  \]
  We do so by chasing the element $\seq[i=1][n]{\id[\avec_i]} \in
  \prod_{i=1}^n\yoneda {\avec_i}{\seq[j=1][1]{\avec_i}}$ around
  $\alpha$'s naturality diagram for the partition
  $\seq[i=1][n]{\seq[j=1][1]{\avec_i} \xfrom{u_i} \partit_i}$ of the
  morphism $u : \avec \xfrom{(u_1, \ldots, u_n)} \bvec$ in $\UU$:
  \diagram{28}
  and so $\inv\Yoneda\compose \Yoneda = \id$.
\end{proof}
We call both this result and the bijection $\Yoneda$ the Yoneda lemma.

\begin{lemma}\emph{Yoneda embedding.}
  The Yoneda multi-functor $\yoneda_{\UU} : \UU \to \Psh\UU$ is
  fully-faithful.
\end{lemma}
\begin{proof}
  Take any $\avec$ and $b$, and consider the Yoneda lemma for the
  representable presheaf $F := \yoneda b$. Take any $f \in \UU(\avec;
  b) = \yoneda b\avec$.  For every partition $\partit$ of every
  $n$-tuple $\cvec$ and $\seq[i=1][n]{u_i : \partit_i \to \avec_i}$,
  we have:
  \[
  (\inv\Yoneda_{\UU, \yoneda b, \avec} f)_{\partit}\seq[i=1][n]{u_i}
  =
  \yoneda b(u_1, \ldots, u_n) f
  =
  f \compose (u_1, \ldots, u_n
  =
  (\yoneda f)_{\partit}\seq[i=1][n]{u_i}
  \]
  and so $\inv\Yoneda_{\UU, \yoneda b, \avec}f = \yoneda_{\UU}f$,
  i.e., $\inv\Yoneda_{\UU, \yoneda b, \avec}$ is the multi-functorial
  action of $\yoneda$ on the multi-homset $\UU(\avec; b)$. Since the
  Yoneda lemma is bijective, this functorial action is bijective,
  i.e., $\yoneda_{\UU}$ is fully-faithful.
\end{proof}
\begin{lemma}\emph{Naturality of the Yoneda Lemma.}
  The Yoneda lemma $\Yoneda : \Psh\UU(\yoneda[\avec], F)
  \xto\isomorphic F\avec$ is natural in the following senses:
  \begin{itemize}
  \item Naturally in $\avec \in \opposite\UU$, i.e.~for all $f : \avec \from \bvec$:
    \[
    Ff (\Yoneda(\yoneda[\avec] \xto\alpha F)) =
    \Yoneda(\yoneda[\bvec] \xto{\yoneda[f]} \yoneda[\avec] \xto\alpha F)
    \]
  \item Multi-naturally in $F \in \Psh\UU$, i.e.~for all $(F_1,
    \ldots, F_n) \xto{\beta} G$ and partition $\partit$ of $\avec$:
    \[
    \beta_{\partit}\seq[i=1][n]{\Yoneda(\yoneda[\partit_i] \xto{\alpha_i} F_i)}
    =
    \Yoneda\parent{\yoneda[\avec] \xto{\beta\compose\seq[i=1][n]{\alpha_i}} G}
    \]
  \end{itemize}
\end{lemma}
We can recover these from the general theory of Yoneda structures, but
we use them to establish the Yoneda structure axioms.
\begin{proof}
  In both cases, employ mindless calculation. For naturality:
  \begin{align*}
    Ff (\Yoneda(\yoneda[\avec] \xto\alpha F))
    &{}=
    Ff \parent{\alpha_{\seq[i=1][n]{\seq[j=1][1]{\avec_i}}}
      \seq[i=1][n]{\id[\avec_i]}}
    \explain={$\alpha$-naturality}
    \alpha_{\seq[i=1][n]{\bvec_{\pair i*}}}
    \seq[i=1][n]{(\yoneda \avec_if_i)\seq[j=1][1]{\id[\avec_i]}}
    \\&{}\explain={identities}
    \alpha_{\seq[i=1][n]{\bvec_{\pair i*}}}
    \seq[i=1][n]{(\yoneda f_i)_{
        \seq[j=1][\abs{\bvec_{\pair i*}}]{\seq[k=1][1]{\bvec_{\pair ij}}}}
      \seq[j=1][\abs{\bvec_{\pair i*}}]{\id[\bvec_{\pair ij}]}}
    \\&{}\explain={$\compose$-def}
    (\alpha\compose\seq[i=1][n]{\yoneda f_i})_{
      \seq[\ell=1][\norm\bvec]{\seq[j=1][1]{\bvec_{\ell}}}}
    \seq[\ell=1][\norm\bvec]{\id[\bvec_\ell]}
    =
    \Yoneda(\alpha \compose\seq[i=1][n]{\yoneda f_i})
  \end{align*}
  as we wanted.

  For multi-naturality:
  \begin{align*}
    \beta_{\partit}\seq[i=1][n]{\Yoneda(\yoneda[\partit_i] \xto{\alpha_i} F_i)}
    &{}=
    \beta_{\partit}\seq[i=1][n]{
      \alpha_{i, \seq[j=1][\abs{\partit_i}]{\seq[k=1][1]{\partit_{\pair ij}}}}
      \seq[j=1][\abs{\partit_{i}}]{\id[\partit_{\pair ij}]}}
    \\&{}\explain={$\compose$-def}
    (\beta\compose\seq[i=1][n]{\alpha_i})_{\partit}
    \seq[\ell=1][\norm \partit]{\id[\partit_{\seq[][-1]{\ell}}]}
    =
    \Yoneda(\beta\compose\seq[i=1][n]{\alpha_i})
  \end{align*}
  as we wanted.
\end{proof}

\begin{lemma}\emph{Special Yoneda reduction.}
  \lemmalabel{special Yoneda reduction}
  For every $\avec \xto{u} b$, we have:
  \(
  \Yoneda_{\UU, \yoneda b, \avec}(\yoneda[\avec] \xto{\yoneda u} \yoneda b)
  = u
  \).
\end{lemma}
\begin{proof}
  Calculate:
  $\Yoneda(\yoneda u) = \yoneda u_{\avec} (\id[\avec_1], \ldots, \id[\avec_n])
  = u \compose (\id, \ldots, \id) = u$.
\end{proof}

\begin{lemma}\lemmalabel{admissible yoneda embedding}
  The Yoneda embedding is, up to equivalence, admissible: for every multi-category $\UU$ there are:
  \begin{itemize}
  \item a multi-category $\Psh'\UU$;
  \item an admissible functor $\yoneda_{\UU}' : \UU \to \Psh'\UU$; and
  \item an adjoint equivalence in $\MultiCAT$, i.e.~an adjunction
    with invertible unit and counit:
    \[
    \quiver{psh adjoint equivalence}
    \]
  \item $2$-cells:
    \begin{center}
      \quiver{psh 2-cell transformer}
      \quiver{psh 2-cell transformer'}
    \end{center}
  \end{itemize}
  such that:
  \begin{center}
    \quiver{psh transformer condition}
    \quiver{psh transformer condition'}
  \end{center}
\end{lemma}
\begin{proof}
  Take any $\avec$ and presheaf $F$. By the Yoneda lemma:
  \[
  \Psh\UU(\yoneda\avec_1, \ldots, \yoneda\avec_n; F) \isomorphic
  F\avec \in \Set
  \]
  is a small set. So now we need to tediously construct a different
  presheaf multi-category $\Psh' \UU$ which is admissible:
  \begin{itemize}
  \item $\Obj{\Psh'\UU} := \Obj{\Psh\UU} + \Obj{\UU}$, i.e., either a
    presheaf or an object, together with a tag telling us which is
    which. We call the left-tagged elements `actual presheaves' and
    the right-tagged elements `formal representables'.  For each $x
    \in \Obj{\Psh'\UU}$, we can associate a presheaf $\sem x$:
    \[
    \sem{\injection_1 F} := F \qquad \sem{\injection_2 a} := \yoneda a
    \]
  \item We define the homs as follows:
    \[
    \Psh'\UU(x_1, \ldots, x_n; y) \definedby \begin{cases}
      \forall i. x_i = \injection_2 a_i: & \sem y(a_1, \ldots, a_n) \\
      \text{otherwise}: & \Psh\UU(\sem{x_1}, \ldots, \sem{x_n}; \sem y)
    \end{cases}
    \]
    We can extend the semantics function into a multi-functor structure
    $\sem- : \Psh'\UU \to \Psh\UU$ using the Yoneda lemma:
    \[
    \sem{f}_{x_1, \ldots, x_n; y} \definedby
    \begin{cases}
      \forall i. x_i = \injection_2 a_i: &
        \inv\Yoneda_{\UU, \sem y, (a_1, \ldots, a_n)}f \\
      \text{otherwise:} & f
    \end{cases}
    \]
  \item The identities are the identity natural multi-transformations
    over actual presheaves, and the identity multi-morphism at $a$
    over the formal representable by $a$. By fiat, the semantics
    functor-structure preserves identities.
  \item To compose multi-morphisms, apply the semantics
    functor-structure to turn them into composable multi-natural
    transformations, compose them in $\Psh\UU$, and then, if needed,
    apply the Yoneda lemma to get the corresponding element of the
    codomain presheaf. By fiat, the semantics functor-structure
    preserves composition.
  \item We therefore also have a functor from the multi-category
    structure $\Psh'\UU$ to $\Psh\UU$.
  \end{itemize}
  The functor $\sem- : \Psh'\UU \to \Psh\UU$ is fully-faithful, and in
  particular faithful, and therefore $\Psh'\UU$ is a multi-category.
  We will take $E := \sem-$, and $E' \definedby \injection_1 : \Psh\UU
  \to \Psh'\UU$.

  We can define the unit and counit, since the component tells us
  whether we need a multi-transformation or merely its
  representation. For the unit we need to split into cases:
  \[
  \unit_{\injection_1F} \definedby \id[F]
  \qquad
  \unit_{\injection_2a} \definedby \id[a] \in \yoneda aa = \sem{\injection_1\sem{\injection_2a}}a
  \]
  For the counit, we know the domain is an actual presheaf, and so use
  the identity natural multi-transformation.

  I expect these are all natural and satisfy the triangle laws.

  The formal Yoneda functor $\yoneda' : \UU \to \Psh'\UU$ sends each
  $a \in \Obj\UU$ to its formal representable and each multi-morphism
  to itself. Since we started with a locally-small category, we have
  by fiat that $\yoneda'$ is admissble:
  \[
  \Psh'\UU(\yoneda'a_1, \ldots, \yoneda'a_n; x) = \sem x(a_1, \ldots,
  a_n) \in \Set
  \]
  The two transformations $\tau : \sem{\yoneda'-} \to \yoneda$ and
  $\tau' : \injection_1\yoneda- \to \injection_1-$ are then
  identities, since the hom definitions at each object take care of
  everything else.

  What makes the two additional laws hold, in the case that we don't
  have an identity, is that the both the unit for the adjunction and
  our case-splitting composition will choose the correct representative.
\end{proof}

\subsection{The Restricted Yoneda multi-functor}
The structure underlying a Yoneda structure is a choice of 4 components:
\begin{itemize}
\item The admissible Yoneda $1$-cells $\yoneda_{\UU} : \UU \to \Psh\UU$,
  together with their presheaf construction. We covered both in the
  previous section, up to equivalence (\lemmaref{admissible
    yoneda embedding}).
\item For every admissible functor $F : \UU \to \UU[2]$, a $1$-cell
  (non-admissible perhaps) $\UU[2](F;-) : \UU[2] \to \Psh\UU$;
\item and a $2$-cell $\map F : \yoneda_{\UU} \to \UU[2](F;F-)$.
\end{itemize}
\begin{center}
$\quiver{yoneda structure structure}$
\end{center}
This section concerns constructing the remainder of this data, and
proving the first axiom of the Yoneda structure.

Let $F : \UU[1] \to \UU[2]$ be an admissible multi-functor. To keep
the notation less ellipsis-cluttered, given a sequence of
objects/multi-morphisms $\xivec$ in $\UU[1]$, we write $F[\xivec]
\definedby \seq[i=1][\abs\xivec]{F\xivec_i}$ for the componentwise
application of $F$.

We define
the structure of the \emph{$F$-restricted Yoneda multi-functor}
$\UU[2](F;-) : \UU[2] \to \Psh\UU$:
\begin{itemize}
  \item It sends each object $x \in \Obj{\UU[2]}$ to the homset functor
    $\UU[2](F[-]; x) : \opposite\UU \to \Set$ sending:
    \begin{itemize}
    \item each sequence $\avec \in \Obj{\opposite\UU}$ to the
      multi-homset $\UU[2](F[\avec]; x)$; and
    \item each morphism $f : \avec \from \bvec$ to the function:
      \[
      \UU[2](F[f]; x) : \parent{(F[\avec]) \xto u x}
      \mapsto (F[\bvec]) \xto{u \compose (F[f])} x
      \]
    \end{itemize}
  \item It sends each morphism $(\xvec) \xto{v} y$, $n \definedby
    \abs\xvec$ to the following natural multi-transformation
    $\UU[2](F; v): \parent{\UU[2](F; \xvec_1), \ldots,
      \UU[2](F;\xvec_n)} \to \UU[2](F;y)$ given, for every $n$-ary
    partition $\partit$ of $\avec$, by:
    \[
    \UU[2](F; v)_{\partit} : \seq[i=1][n]{F[\partit_i]\xto{v_i} \xvec_i}
    \mapsto (F[\avec])\xto{v\compose(v_1, \ldots, v_n)} y
    \]
\end{itemize}

\begin{lemma}
  These data define a multi-functor $\UU[2](F;-) : \UU[2] \to
  \Psh\UU$.
\end{lemma}
\begin{proof}
  The restricted Yoneda functor is well-defined:
  \begin{itemize}
  \item Each hom $\UU[2](F[\avec]; x)$ is a set since $F$ is admissible.
  \item Given a morphism $\avec \xfrom f \bvec$ in $\UU$, $F[\avec]
    \xfrom{F[f]} F[\bvec]$ in $\UU[2]$, and so given $F[\avec] \xto{u}
    x$ the composition $u \compose (F[f]) : F[\bvec] \to x$
    type-checks as a well-defined multi-morphism.
  \item For identity preservation, calculate:
    \[
    \UU[2](F[{\id[\avec]}]; x)(u) \definedby
      u \compose\seq[i=1][\abs\avec]{F \id[ \avec_i]}
    \explain={$F$ multi-functor}
      u \compose\seq[i=1][\abs\avec]{  \id[F\avec_i]}
    \explain={multi-category}
      u
    \]
    and so $\UU[2](F[{\id[\avec]}]; x)  = \id[{\UU[2](F[\avec]; x)}]$.
  \item Take two composable morphisms in $\opposite\UU$, i.e.: $\avec
    \xfrom f \bvec \xfrom g \cvec$ in $\UU$ and calculate for a
    generic $u : (F[\avec]) \to x$:
    \begin{align*}
      \UU[2](F[f \compose g]; x)u
      &{}= u \compose (F[f \compose g])
      \explain={$\compose$-def in $\opposite\UU$}
      u \compose \seq[i=1][\abs f]{F\parent{f_i \compose
        \seq[j=1][\abs {g_{\pair i*}}]{g_{\pair ij}}}}
      \\&{}\explain={$F$ multi-functor}
      u \compose \seq[i=1][\abs f]{Ff_i\compose
        \seq[j=1][\abs {g_{\pair i*}}]{Fg_{\pair ij}}}
        \explain={associativity of composition\\in a multi-category}
      (u \compose \seq[i=1][\abs f]{Ff_i})\compose
        \seq[\ell=1][\abs g]{Fg_\ell}
      \\{}&=
      \UU[2](F[g]; x) \compose \UU[2](F[f]; x) u
    \end{align*}
    and so $\UU[2](F[f \compose g]; x) = \UU[2](F[g]; x) \compose \UU[2](F[f]; x)$.
    Thus, $\UU[2](F;x) : \opposite\UU \to \Set$ is a presheaf.
  \end{itemize}

  To show that the multi-functorial action is well-defined, take
  $(\xvec) \xto v y$ in $\UU[2]$ with $n \definedby \abs\xvec$, and
  show the naturality condition by taking an $n$-ary partition
  $\seq[i=1][n]{\partit_i \xfrom{\phi_i} \partit'_i}$ of some $\avec
  \xfrom f \bvec$ in $\UU$, chasing an element of
  $\prod_{i=1}^n\UU[2](F[\avec];\xvec _i)$:
  \diagram{29}

  To show this action is multi-functorial, take an identity $\id[x]$
  in $\UU[2]$ and the unique unary partition $\seq{\avec}$ of some
  $\avec$:
  \diagram{30}
%
  Next, given composable multi-morphisms $(\yvec) \xto v z$, $n
  \definedby \abs\yvec$ and $\seq[i=1][n]{(\xvec_i) \xto v_i \yvec_i}$
  in $\UU[2]$, letting $m \definedby \norm\xvec$ and $m_i \definedby
  \abs{\xvec_i}$ and taking any $m$-ary partition of some $\avec$,
  calculate: \diagram{31}
  So $\UU[2](F;v\compose\seq[i=1][n]{v_i})=
  \UU[2](F;v)\compose\seq[i=1][n]{\UU[2](F;v_i)}$ and we have a
  functor.
\end{proof}

The multi-natural transformation $\map F : \yoneda_{\UU} \To
\UU[2](F;F-)$ we now define internalises the multi-functorial action of
$F$. Explicitly, the component at $b \in \Obj\UU$ is given by the
natural multi-transformation $(\map F)_{b} : (\yoneda_{\UU}b) \to
\UU[2](F;Fb)$ whose component at any unary partition
$\seq[i=1][1]\avec$ of some sequence $\avec$ is given by:
\[
(\map F)_{b, \seq[i=1][1]\avec} :
  \prod\set{\yoneda b\avec} \ni
  \seq[i=1][1]{(\avec) \xto{u} b} \mapsto
  (F[\avec] \xto{Fu} Fb) \in \UU[2](F[\avec]; Fb)
\]
or, equivalently, using the Yoneda lemma:
\[
(\map F)_{b, \seq[i=1][1]\avec} \seq{u} \definedby \Yoneda_{F;\avec}(\yoneda(F[\avec])\xto{\yoneda (Fu)}\yoneda (Fb)) \in \yoneda (Fb)(F[\avec]) = \UU[2](F[\avec];Fb)
\]

\begin{lemma}
  This data defines a multi-natural transformation:
  $\map F : \yoneda_{\UU} \To \UU[2](F;F-)$.
\end{lemma}
\begin{proof}
  By Yoneda reduction, the two definitions are
  equivalent. We need to show:
  \begin{itemize}
  \item Well-definedness: each component is multi-natural.
  \item Naturality of the components over-all.
  \end{itemize}
  Unfortunately, these don't just follow automatically from our
  formulation using the Yoneda lemma together with its naturality and
  multi-naturality, since we're applying $F$ to $u$. However, the
  formulation using Yoneda might give some intuition why this
  operation is multi-natural, and our later results will allow us to
  recover it internally.

  To show well-definedness, we need to show that each component $(\map
  F)_{b}$ is a natural multi-transformation. Take the unique unary
  partition $\seq[i=1][1] f$ of some morphism $f : \avec \from \avec'$
  in $\UU$, and chase an element of $\prod\set{\yoneda b\avec}$ around the
  naturality square:
  \diagram{32}

  To show multi-naturality, take any $f : \bvec \to b$ in $\UU$, let
  $n \definedby \abs\bvec$ and consider one component of the two sides
  of the multi-naturality diagram in $\Psh\UU$:
  \[
  (\map F)_{b}\compose(\yoneda f)
  =
  \UU[2](F;Ff)\compose\seq[i=1][n]{(\map F)_{\bvec_i}}
  \]
  by taking an $n$-ary partition $\partit$ of some $\avec$ and
  calculating:
  \diagram{33}
  and so $\map F$ is multi-natural.
\end{proof}

\begin{proposition}\emph{Axiom 1.}
  The multi-functor $\UU[2](F;-)$ together with the multi-natural
  transformation $\map F$ form the left Kan extension of
  $\yoneda_{\UU}$ along $F$:
  \begin{center}
    $\quiver{yoneda structure structure}$
  \end{center}
\end{proposition}
\begin{proof}
  Like the Yoneda lemma, we roll up our sleeves and just do
  it. We need to show a bijective correspondence:
  \begin{center}
    $\quiver{axiom 1 spelled out}$
  \end{center}

  Starting with uniqueness, let $\alpha := (\beta \hcompose F)
  \vcompose (\map F)$ be the pasted multi-natural transformation.


  First, unpack the determining structure of $\beta$.  Take any
  component $x \in \Obj\UU[2]$ of the multi-natural transformation
  $\beta : \UU[2](F;-) \To G$, giving a natural multi-transformation
  $\beta_{x} : (\UU[2](F;x)) \to Gx$. Take any component of
  $\beta_{x}$, i.e.~a uniary partition of some $\avec \in
  \Obj{\opposite\UU}$ and we now have:
  \[
  \beta_{x, \seq[k=1][1]{\avec}} :
  \UU[2](F[\avec]; x) \to Gx(\avec)
  \]
  and we need to determine the image of $\beta_{x, \seq[k=1][1]\avec}$
  on some $F[\avec]\xto v x$, and show that it only depends on $\alpha$.

  Note that we have a natural multi-transformation $Gv : (GF[\avec])
  \to Gx$ in $\Psh\UU$ and so we may use the naturality of each
  component of $\alpha$, form the composite and apply the Yoneda
  lemma to it:
  \[
  Gv \compose \seq[i=1][n]{\yoneda \avec_i \xto{\alpha_{\avec_i}} GF\avec_i}
  : (\yoneda[\avec]) \to Gx
  \implies
  \Yoneda_{Gx, \avec}(Gv\compose \seq[i=1][n]{\alpha_{\avec_i}}) \in Gx(\avec)
  \]
  We will show that:
  \[
  \beta_{x, \seq[k=1][1]{\avec}}\seq[k=1][1]v =
    \Yoneda_{Gx, \avec}(Gv\compose \seq[i=1][n]{\alpha_{\avec_i}})
  \]
  and since the latter only involves $\alpha$, and not $\beta$, then
  pasting with $\map F$ is injective. Indeed:
  \begin{align*}
    \Yoneda_{Gx, \avec}(Gv\compose \seq[i=1][n]{\alpha_{\avec_i}})
    &{}=
    (Gv\compose \seq[i=1][n]{\alpha_{\avec_i}})
      _{\seq[i=1][n]{\seq[j=1][1]{\avec_i}}}\seq[i=1][n]{\id[\avec_i]}
    \\&{}=
    (Gv)_{\seq[i=1][n]{\seq[j=1][1]{\avec_i}}}\seq[i=1][n]{\alpha_{\avec_i}(\id[\avec_i])}
    \\&{}\explain={$\alpha = (\beta \hcompose F)\vcompose (\map F)$}
    (Gv)_{\seq[i=1][n]{\seq[j=1][1]{\avec_i}}}\seq[i=1][n]{\beta_{F\avec_i}(F\id[\avec_i])}
    \\&{}\explain={multi-functoriality}
    (Gv)_{\seq[i=1][n]{\seq[j=1][1]{\avec_i}}}\seq[i=1][n]{\beta_{F\avec_i}(\id[F\avec_i])}
    \\&{}\explain={$\compose$-def in $\Psh\UU$}
    (Gv\compose \seq[i=1][n]{\beta_{F\avec_i}})
       _{\seq[i=1][n]{\seq[j=1][1]{\avec_i}}}\seq[i=1][n]{\id[F\avec_i]}
    \\&{}\explain={multi-naturality}
    (\beta_x \compose (\UU[2](F;v)))
      _{\seq[i=1][n]{\seq[j=1][1]{\avec_i}}}\seq[i=1][n]{\id[F\avec_i]}
    \\&{}\explain={$\compose$-def}
    \beta_{x, \seq[k=1][1]{\avec}} \seq[k=1][1]{
      \UU[2](F;v)_{\seq[i=1][n]{\avec_i}}
      \seq[i=1][n]{\id[F\avec_i]}}
    \\&{}\explain={$\UU[2](F;v)$-def}
    \beta_{x, \seq[k=1][1]{\avec}} \seq[k=1][1]{
      v\compose\seq[i=1][n]{\id[F\avec_i]}}
    \explain={multi-category}
    \beta_{x, \seq[k=1][1]{\avec}} \seq[k=1][1]{v}
  \end{align*}
  and so the correspondence is injective.

  To show surjectivity, consider any multi-natural $\alpha$, and set:
  \[
  \beta_{x, \seq[k=1][1]{\avec}}\seq[k=1][1]v \definedby
    \Yoneda_{Gx, \avec}(Gv\compose \seq[i=1][n]{\alpha_{\avec_i}})
  \]
  We need to show:
  \begin{itemize}
  \item $\beta_{x}$ is a natural multi-transformation;
  \item $\beta$ is a multi-natural transformation;
  \item pasting with $\map F$ recovers $\alpha$, i.e.: $\alpha = (\beta
    \hcompose F)\vcompose (\map F)$.
  \end{itemize}

  To show that $\beta_{x} : \UU[2](F;x) \to Gx$ is natural, take any
  $f : \avec \from \bvec$ in $\UU$ and calculate:
  \begin{align*}
  Gx f(\beta_{x, \seq[k=1][1]{\avec}}\seq[k=1][1]v)
  &{}=
  Gx f\parent{\Yoneda(Gv \compose \seq[i=1][\abs\avec]{\alpha_{\avec_i}})}
  \explain={$\Yoneda$-naturality}
  \Yoneda\parent{Gv \compose
    \seq[i=1][\abs\avec]{\alpha_{\avec_i}\compose \yoneda f_i}}
  \\&{}\explain={$\alpha$-multi-naturality}
  \Yoneda\parent{Gv \compose
    \seq[i=1][\abs\avec]{GF f_i\compose
      \seq[j=1][\abs{\bvec_{\pair i*}}]{
        \alpha_{\bvec_{\pair ij}}}}}
  \\&{}\explain={$G$-multi-functor}
    \Yoneda\parent{G(v \compose Ff_i)\compose
      \seq[\ell=1][\abs{\bvec}]{
        \alpha_{\bvec_{\ell}}}}
  =\beta_{x, \seq[k=1][1]{\bvec}}\seq[k=1][1]{\UU[2](F[f];x)v}
  \end{align*}

  To show that $\beta$ is multi-natural, take any $v : (\xvec) \to y$
  in $\UU[2]$ and $v_i : F[\avec_i] \to \xvec_i$ for $1 \leq i \leq n
  \definedby \abs\avec$:
  \begin{align*}
    (Gv \compose \seq[i=1][n]{\beta_{\xvec_i}})
      _{\seq[i=1][n]{\avec_i}}(v_1, \ldots, v_n)
    &{}=
    Gv\seq[i=1][n]{\Yoneda(Gv_i \compose
      \seq[j=1][\abs{\avec_i}]{\alpha_{\avec_{\pair ij}}})}
    \\&\explain={$\Yoneda$-multi-naturality, associativity}
    \Yoneda((Gv\compose\seq[i=1][n]{Gv_i}) \compose
    \seq[\ell=1][\norm{\avec}]{\alpha_{\avec_{\ell}}})
    \\&{}\explain={$G$-multi-functoriality}
    \Yoneda(G(v\compose\seq[i=1][n]{v_i}) \compose
    \seq[\ell=1][\norm{\avec}]{\alpha_{\avec_{\ell}}})
    \\&{}=
    (\beta_{y}\compose \UU[2](F;v))\seq[\ell=1][n]{v_1, \ldots, v_n}
  \end{align*}

  To show that we recover $\alpha$, take any $u : (\avec) \to b$ and
  calculate:
  \begin{align*}
    ((\beta \hcompose F) \vcompose (\map F))_{b, \seq[k=1][1]\avec}(u)
    =
    \beta(Fu)
    =
    \Yoneda(GFu \compose \seq[i=1][\abs\avec]{\alpha_{\avec_i}})
    \explain={$\alpha$-multi-naturality}
    \Yoneda(\alpha_{b}\compose\yoneda u)
    =\alpha_{b, \seq[k=1][1]\avec}u
  \end{align*}
  concluding the proof.
\end{proof}

\subsection{The Relative Yoneda lemma}
\begin{proposition}\emph{Axiom 2.}
  For every multi-functor (admissible or not) $G : \UU[3] \to \UU$,
  the following triangle is a left Kan lifting of $\yoneda_{\UU}$
  through $\UU[2](F;-)$:
  \begin{center}
    $\quiver{axiom-2 lifting}$
  \end{center}
  I.e., for every $H : \UU[3] \to \UU[2]$, pasting
  $(\map F)\hcompose G$ along $\UU[3] \xto G \UU \xto F \UU[2]$ gives a bijection:
  \begin{center}
    $\quiver{axiom-2 bijection}$
  \end{center}
  We call the bijection $\alpha \mapsto \Yoneda_{\UU,
    F,H,G\seq-}\alpha$ the \emph{Yoneda lemma relative to $G$ along
    $F$}.
\end{proposition}
\begin{proof}
  Starting with uniqueness, given the data:
  \begin{center}
    $\quiver{axiom-2 proof: uniqueness data}$
  \end{center}
  such that
  \begin{center}
    $\quiver{axiom-2 proof: uniqueness assumption}$
  \end{center}
  we will show that $\beta$ is uniquely determined by $\alpha$,
  i.e., the relative Yoneda lemma is injective.

  Take any $e \in \Obj{\UU[3]}$, and then: $\beta_e : (FGe) \to He$ is
  a multi-morphism in $\UU[2]$, i.e.~an element $\beta_e \in
  \UU[2](F[(Ge)]; He)$, and so by the Yoneda lemma we have a
  multi-natural transformation:
  \[
  \alpha'_e \definedby
  \inv\Yoneda_{\UU,\UU[2](F;He),\seq{Ge}}\beta_e : \yoneda (Ge) \to
  \UU[2](F[-];He)
  \]
  which is the same type as the component $\alpha_e$. We will show
  that $\alpha'_e = \alpha_e$, and so $\beta_e$ is uniquely determined
  by $\alpha$:
  \[
  \beta_e = \Yoneda_{\UU,\UU[2](F;He),\seq{Ge}}\alpha_e
  \]
  Take any $\avec \in \Obj{\opposite\UU}$ and $u : (\avec) \to Ge$ in
  $\UU$. On the one hand:
  \begin{align*}
    \alpha'_{e, \avec}(u)
    &{}\explain={$\Yoneda$-def}
    \UU[2](F[u];He)(\beta_e)
    =
    \beta_e\compose(Fu)
  \end{align*}
  On the other hand:
  \begin{align*}
    \alpha_{e, \avec}(u)
    &{}\explain={assumption}
    \parent{(\UU[2](F;-)\hcompose\beta)\vcompose
      (\map F \hcompose G)}_{e,\avec}(u)
    \explain={$(\hcompose$-def$)\vcompose(\map F$-def$)$}
    \UU[2](F;\beta_e)(Fu)
    =
    \beta_e \compose (Fu)
  \end{align*}
  and therefore $\beta$ is uniquely determined by $\alpha$.

  For the existence, set $\beta_e \definedby
  \Yoneda_{\UU,\UU[2](F;He),\seq{Ge}}\alpha_e$ --- this choice
  motivates our notation --- and we'll show:
  \begin{itemize}
  \item $\beta : FG \to H$ is multi-natural; and
  \item $\alpha = (\UU[2](F;-)\hcompose\beta)\vcompose (\map F
    \hcompose G)$
  \end{itemize}
  For multi-naturality, take any $(e_1, \ldots, e_n) \xto{w} e$ in
  $\UU[3]$. We want to prove:
  \[
  Hw \compose (\beta_{e_1}, \ldots, \beta_{e_n}) = \beta_e \compose (FGw)
  \]
  Calculate:
  \begin{align*}
    Hw \compose (\beta_{e_1}, \ldots, \beta_{e_n})
    &{}=
    \UU[2](F;Hw)\seq[i=1][n]{\Yoneda(\alpha_{e_i})}
    \explain={$\Yoneda$-multi-naturality}
    \Yoneda\parent{\UU[2](F;Hw)\compose\seq[i=1][n]{\alpha_{e_i}}}
    \\&{}\explain*={$\alpha : \yoneda(G-) \to \UU[2](F;H-)$\\multi-naturality}
    \Yoneda\parent{\alpha_{e}\compose(\yoneda (Gw))}
    \explain*={$\Yoneda$-naturality}
    \UU[2](F[Gw]; He)(\Yoneda(\alpha_e))
    =
    \beta_e \compose (FGw)
  \end{align*}
  as we wanted.

  To show that we recover $\alpha$ by pasting, take any $e \in
  \UU[3]$, $\avec \in \Obj{\opposite\UU}$, $u : (\avec) \to Ge$. We
  need to show:
  \[
  \alpha_{e, \avec}(u) = \UU[2](F;\beta_e)(\map F u) = \beta_e \compose (F u)
  \]
  Calculate:
  \begin{align*}
    \beta_e \compose (F u)
    &{}=
    \UU[2](F[u]; He)(\Yoneda(\alpha_e))
    \explain={$\Yoneda$-naturality}
    \Yoneda\parent{
      \alpha_e \compose (\yoneda u)
    }
    \explain={$\Yoneda$-multi-naturality}
    \alpha_{e,\avec}(\Yoneda(\yoneda u))
    \explain={Yoneda reduction%
      % vertical alignment hack --- this looks nicer
      \vphantom{y}}
    \alpha_{e,\avec}u
  \end{align*}
  as we wanted.
\end{proof}

\subsection{The presheaf-construction pseudo-functor}
The third axioms concerns internalising the presheaf construction into
a pseudo-functor from the locally-small multi-categories to
multi-categories. Doing so uses an internalisation of the hom
profunctor that extends the postulate structure $\UU[2](F;-)$, first
to an action on $F$, and then to a similar action in the covariant
position.

\paragraph{Internalising the hom profunctor}
The following structure and properties follow from Axiom~1, letting
$\AdMultiCAT$ denote that full sub-$2$-category of $\AdMultiCAT$
consisting of the admissible multi-functors, we have, for every triple
of multi-categories $\UU[1], \UU[2], \UU[3]$, with $\UU[1]$
locally-small, a functor:
\[
\UU[2](-,- -) :
\opposite{\AdMultiCAT(\UU[1], \UU[2])}\times \MultiCAT(\UU[3], \UU[2]) \to
\MultiCAT(\UU[3], \Psh{\UU[1]})
\]
To define it, we first define its partial application with $\UU[3] :=
\UU[2]$, $\Id : \UU[3] \to \UU[2]$, namely:
\[
\UU[2](-,\Id-) :
\opposite{\AdMultiCAT(\UU[1], \UU[2])}\to
\MultiCAT(\UU[2], \Psh{\UU[1]})
\]
given by:
\begin{itemize}
\item sending `objects', i.e., admissible multi-functors $F : \UU[1] \to
  \UU[2]$, to the postulated multi-functor:
  \[
  \UU[2](F,\Id-) \definedby \UU[2](F;-) : \UU[2] \to \Psh{\UU[1]}
  \]
\item sending `morphisms', i.e., multi-natural transformations
  $\quiver{hom-profunctor morphism input type}$, to the unique
  multi-natural transformation:
  \begin{center}
  $\quiver{hom-profunctor morphism output type}$
  \end{center}
  satisfying:
  \begin{center}
    $\quiver{hom-profunctor special case up}$
  \end{center}
  Working this action out explicitly, it's given by the multi-natural
  transformation consisting of natural multi-transformations acting by
  pre-composition with the appropriate component of $\alpha$:
  \[
  \UU[2](\alpha;\Id-)_{x,\avec} :
    \parent{F[\avec] \xto u x} \mapsto
    \parent{G[\avec] \xto{\seq[i]{\alpha_{\avec}}} F[\avec] \xto u}
  \]
\end{itemize}
We then define the more general functor using horizontal composition
as follows. Given:
\begin{center}
  \quiver{hom profunctor general input}
\end{center}
we define:
\begin{center}
  $\quiver{hom profunctor general def via hcompose}$
\end{center}

Explicitly, the functor $\UU[2](-;- -)$ is given by:
\begin{itemize}
\item sending `objects', i.e., admissible multi-functors $F : \UU[1]
  \to \UU[2]$ and arbitrary multifunctor $H : \UU[3] \to \UU[2]$, to
  the composed multi-functor:
  \[
  \UU[2](F,H-) \definedby \UU[2](F;-)\compose H : \UU[3] \to \UU[2]
  \to \Psh{\UU[1]}
  \]
  Its object function is given by:
  \begin{itemize}
  \item for all $x \in \Obj{\UU[2]}$ and $\avec\in \Obj{\opposite\UU}$:
    \[
     \UU[2](F;H-)x\,\avec \definedby \UU[2](F[\avec];Hx)
    \]
  \item for all $x \in \Obj{\UU[2]}$ and $\avec \xfrom{f} \bvec$ in $\UU$:
    \[
    \UU[2](F;H-)x\,f \definedby \UU[2](F[f];Hx) : ((F[\avec]) \xto u Hx) \mapsto
      \parent{(F[\bvec]) \xto{F[f]} (F[\avec]) \xto u Hx}
    \]
  \end{itemize}
  Its hom action is given by:
  \begin{itemize}
\item for all $g : \xvec \to y$ in $\UU[2]$ and $\avec \in \Obj{\opposite\UU}$:
  \[
  \UU[2](F;Hg)_{\avec} \definedby
  \UU[2](F[\avec];Hg) : \seq[i=1][\abs\xvec]{(F[\avec]) \xto{u_i} Hx} \mapsto
  \parent{F[\avec]\xto{Hg\compose \seq[i]{u_i}} Hy}
  \]
  \end{itemize}
\item and its operation on `morphisms':
  \begin{center}
  \quiver{hom profunctor general input-output}
  \end{center}
  is given by the multi-natural transformation whose component at $x$ is the
  natural multi-transformation whose component at $\avec$ is:
  \[
  \UU[2](\alpha;\gamma)_{x, \avec} : \parent{(F[\avec]) \xto{u} Hx} \mapsto
  \parent{(G[\avec]) \xto{\seq[i=1][\abs\avec]{\alpha_{\avec_i}}}
      (F[\avec]) \xto{u} Hx \xto{\gamma_x} Kx}
  \]
\end{itemize}

Street and Walters then use this functor, still only with Axiom~1, to
extend the presheaf construction to a co-lax (2-)functor from the
$2$-category of locally-small multi-categories to multi-categories
$\Psh- : \coop{\LocallySmall\MultiCAT} \to \MultiCAT$, sending:
\begin{itemize}
\item each multi-category $\UU$ to its presheaf construction $\Psh\UU$;
\item each (necessarily admissible) multi-functor $F : \UU \from \UU[2]$
  between locally-small categories to the multi-functor:
  \(
  \Psh F \definedby \Psh{\UU[1]}(\UU[1](\Id;F-), -): \Psh \UU \to \Psh{\UU[2]}
  \) given explicitly by:
  \begin{itemize}
  \item sending each presheaf $P \in \Psh\UU$ to the presheaf $\Psh F
    P := \Psh{\UU[1]}(\UU[1](\Id;F-), P) \in \Psh{\UU[2]}$:
    \begin{itemize}
    \item sending each $\xvec \in \opposite{\UU[2]}$ to the set of
      $\abs\xvec$-ary natural multi-transformations of type $\beta :
      (\UU[1](\Id;F[\xvec])) \to P$, i.e. $\beta : \yoneda_{\UU[1]}(F[\xvec]) \to P$;
    \item sending each $\xvec \xfrom{g} \yvec$ in $\UU[2]$ to the
      natural multi-transformation that acts by pre-composing with the
      multi-transformation that acts by post-composing $Fg$:
      \begin{align*}
      \Psh F g : \parent{\yoneda_{\UU[1]}(F[\xvec]) \xto\beta P}
      &{}\mapsto
      \parent{\yoneda_{\UU[1]}(F[\yvec]) \xto{\yoneda_{\UU[1]}(F[g])}
        \yoneda_{\UU[1]}(F[\xvec])
        \xto\beta P}
      \\&{}=
      \seq[\partit]{
        \seq[\ell=1][\abs\yvec]{(\partit_\ell)\xto{u_\ell} F\yvec_i}
        \mapsto
        \beta_{\partit_{\pair *{\++}}}\compose(Fg_i\compose\seq[j=1][\abs{\yvec_{\pair i*}}]{u_{\pair ij}})
      }
      \end{align*}
    \end{itemize}
  \item sending each natural multi-transformation $\alpha : (\Pvec)
    \to Q$ in $\Psh\UU$ to the natural multi-transformation $\Psh
    F\alpha : (\Psh F [\Pvec]) \to \Psh F Q$ in $\Psh{\UU[2]}$ given
    at partition $\partit$ of $\xvec \in \Obj{\UU[2]}$ by:
    \[
    \seq[i=1][\abs\Pvec]{\yoneda(F[\partit_i]) \xto{\beta_i} P_i}
    \mapsto
    \parent{
      \seq[\ell=1][\abs\xvec]{\yoneda (F\xvec_\ell)}
      \xto{
        \alpha_{F[[\partit]]}\compose\seq[i=1][\abs\Pvec]{\beta_{i, \seq[j=1][1]{\partit_i}}}
      }
      Q\xvec
    }
    \]
  \end{itemize}
\item each multi-natural transformation $\quiver{multi-natural
  transformation coop}$ to the multi-natural transformation whose
  $\quiver{psh 2-cell action type}$ component at $P \in \Psh\UU$ is
  the unary natural multi-transformation whose component at a
  partition $\seq[j=1][1]{\xvec}$ of some $\xvec \in
  \opposite{\UU[2]}$ is:
  \[
  \Psh\alpha_{P,\partit} :
    \seq[j=1][1]{\parent{\yoneda(F[\xvec])} \xto{\beta} P}
    \mapsto
    \parent{\parent{\yoneda(F[\xvec])}
      \xto{\beta \compose \seq[i = 1][\abs\xvec]{\yoneda\alpha_{\xvec_i}}} P}
    \]
\end{itemize}
To exhibit $\Psh- : \coop{\LocallySmall\MultiCAT} \to \MultiCAT$ as a
co-lax $2$-functor, we need two families $2$-cells:
\begin{itemize}
\item For every locally-small multi-category $\UU[1]$, a multi-natural
  transformation:
  \begin{center}
  $\quiver{iota type}$
  \end{center}
  It is given by factorising the identity multi-natural transformation
  between the Yoneda embedding and itself:
  \begin{center}
    $\quiver{internal yoneda lemma}$
  \end{center}
  We denote it by $\Yoneda$ because it is given by the Yoneda
  lemma. We can see that directly by calculation (the factorisation
  maps $\id$ to $(\yoneda \xto u F) \mapsto \Yoneda (\Id u \compose
  \id) = \Yoneda u$), but also by noting:
\begin{itemize}
\item The components $\Yoneda_{F,\avec}$ are natural in $\avec$ and so
  we have a natural multi-transformation $\Yoneda_{F}:
  \Psh\UU(\yoneda; F) \to F$, and these are further multi-natural in
  $F$ and so we have a multi-natural transformation $\Yoneda :
  \Psh\UU(\yoneda;-) \to \Id$.
\item By Yoneda reduction:
  \[
  ((\Yoneda \hcompose \yoneda) \vcompose (\map \yoneda)) (u)
  = \Yoneda(\yoneda u) = u
  \]
  and so $\Yoneda$ is the left extension.
\end{itemize}
\item We will also need, for every triple of locally-small
  multi-catgeories and composable multi-functors $\UU[1] \xfrom F
  \UU[2] \xfrom G \UU[3]$ a $2$-cell:
  \begin{center}
    $\quiver{gamma type}$
  \end{center}

  Street and Walters axiomatise a more general structure with a
  stronger axiom. For a composable pair of multi-functors $\UU[1]
  \xfrom H \UU[2] \xfrom G \UU[3]$ consisting of a (necessarily
  admissible) multi-functor $G$ between locally-small multi-categories
  and admissible functor $K$, a $2$-cell:
  \begin{center}
    $\quiver{theta type}$
  \end{center}
  This $2$-cell specialises to $\gamma_{F,G}$ taking $\UU[1]
  \definedby \Psh\UU$ and
  \[
  H \definedby \Psh\UU(-,(\yoneda \compose F)-) :
  \Psh{\UU[1]} \xfrom{\yoneda} \UU[1] \xfrom F \UU[2]
  \]
  The defining universal property for $\theta_{G, H}$ is that it's
  the unique factorisation:
  \begin{center}
    \quiver{theta UP}
  \end{center}
  It consists, for each $\UU[1]$ object $a$, and
  $\opposite{\UU[3]}$ object $\zvec$, of the functions:
  \[
  \theta_{G, H, a, \seq[j=1][1]{\zvec}} : \seq[j=1][1]{HG[\zvec] \xto{u} a}
    \mapsto
    \seq[\yvec \in \opposite{\UU[2]}]
        { \seq[i=1][\abs\zvec]{\yvec_i \xto{v_i} G[\zvec_i]}
          \mapsto
          \parent{H[\yvec]\xto{u\compose\seq[i=1][\abs\zvec]{Hv_i}} a}
    }
  \]
\end{itemize}

Axiom~3 states that the $2$-cells $\Yoneda$ and $\theta$ are
$2$-isomorphisms, and in particular making $\Psh- :
\coop{\LocallySmall\MultiCAT} \to \MultiCAT$ a pseudo-functor.  Before
we prove it, we'll first characterise the multi-natural isomorphisms,
and in presheaf multi-categories that will require us to characterise
the natural multi-isomorphisms.

\begin{lemma}\emph{Isomorphism lemma.}
  \begin{itemize}
  \item
    Let $\UU$ be a multi-category, $P, Q \in \Psh\UU$ presheaves over
    $\UU$, and $\alpha : (P) \to Q$ a \emph{unary} natural
    multi-transformation. Then $\alpha$ is an isomorphism iff all of
    its components are bijective.
  \item
    A multi-natural transformation $\quiver{multi-natural
      transformation type}$ is an isomorphism (w.r.t. $2$-cell
    vertical composition) iff its components $\alpha_{a} : (Fa) \to
    Ga$ are isomorphisms in $\UU[2]$.
  \end{itemize}
\end{lemma}
\begin{proof}
  The `only if' directions follow since the identities are
  component-wise identity functions/morphisms and the appropriate
  composition is given component-wise (up-to rearrangement of tuples
  via a canonical isomorphism). So the remainder of proof amounts to
  showing that appropriate multi-naturality and naturality conditions
  for the assembled component-wise inverse.

  Starting with a component-wise bijective natural
  multi-transformation $\alpha : (P) \to Q$, define its inverse
  componentwise, at each unary partition $\seq[j=1][1]{\avec}$:
  \[
  (\inv\alpha)_{\seq[j=1][1]{\avec}} : \prod\set{Q\avec} \xto\isomorphic Q\avec
  \xto{\inv{\parent{\alpha_{\seq[j=1][1]{\avec}}}}}
  \prod\set{P\avec}
  \xto\isomorphic
  P\avec
  \]
  A straightforward calculation shows that this multi-transformation
  would be the inverse multi-morphism, provided it is indeed
  natural. Take any $f : \avec \from \bvec$ in $\UU$, and calculate:
  \diagram{34}

  Next, consider a multi-natural transformation $\quiver{multi-natural
    transformation type}$ whose components are
  multi-isomorphisms. Define $(\inv\alpha)_{a} : (Ga) \to Fa$ by:
  \[
  (\inv\alpha)_{a} \definedby \inv{\parent{\alpha_{a}}}
  \]
  and we now need to show $\inv\alpha : G \to F$ is
  multi-natural. Take any $f : (\avec) \to b$ in $\UU[1]$, let $n
  \definedby \abs\avec$, and calculate: \begin{align*}
    Gf \compose
    \seq{(\inv\alpha)_{\avec_i}}
    &{}\explain={identity + $\inv\alpha$-def}
    \id[b]\compose (Gf \compose
    \seq[i=1][n]{\inv{\parent{\alpha_{\avec_i}}}})
    \explain={inverses}
    \parent{\inv{\parent{\alpha_{b}}}
      \compose (\alpha_{b})}
    \compose (Gf \compose
    \seq[i=1][n]{\inv{\parent{\alpha_{\avec_i}}}})
    \\&{}\explain={$\compose$-assoc.}
    \inv{\parent{\alpha_{b}}}
      \compose \parent{\alpha_{b}
    \compose (Gf \compose
    \seq[i=1][n]{\inv{\parent{\alpha_{\avec_i}}}})}
    \explain={multi-naturality}
    \inv{\parent{\alpha_{b}}}
      \compose \parent{Ff \compose
        \seq[i=1][n]{\alpha_{\avec_i}\compose (\inv{\parent{\alpha_{\avec_i}}})}}
    \\&{}\explain={inverses}
    \inv{\parent{\alpha_{b}}}
      \compose \parent{Ff \compose
        (\seq[i=1][n]{\id[\avec_i]})}
    \\&{}\explain={identities}
    \inv{\parent{\alpha_{b}}}
      \compose (Ff)
  \end{align*}
  as we want.  This proof is the usual one presented algebraically
  since we don't have a convenient diagrammatic formalism in the
  multi-categorical setting.
\end{proof}

\begin{proposition}\emph{Axiom 3.}
  \begin{itemize}
  \item
    The multi-natural transformation $\Yoneda : \Psh{\Id[\UU]} \to
      \Id[\Psh\UU]$ is a multi-natural isomorphism.
  \item
    The multi-natural transformations $\theta_{G, H} : \UU[3](H
    \compose G, -) \to \Psh G \compose \UU(H, -)$ are multi-natural
    isomorphisms.
  \end{itemize}
\end{proposition}
\begin{proof}
  By the Yoneda lemma, each component of $\Yoneda$ is a bijection, and
  by the Isomorphism Lemma (1) each component $\Yoneda_P$ is a natural
  multi-isomorphism, and therefore (2) $\Yoneda$ is a multi-natural
  isomorphism.

  Turning to $\theta_{G, H}$ for $\UU[1] \xfrom H \UU[2] \xfrom G
  \UU[3]$, we will use its UP to show that each component:
  \[
  \theta_{G, H, a,
    \seq[k=1][1]\zvec} : \UU[1](HG[\zvec], a) \to
  \Psh{\UU[2]}(\yoneda[G[\zvec]], \UU[1](H, a))
  \]
  is in fact the inverse Yoneda Lemma:
  \[
  \inv\Yoneda_{\UU[1]; \UU(H[-],a), G[\zvec]} : \UU[1](HG[\zvec], a) \to
  \Psh{\UU[2]}(\yoneda[G[\zvec]], \UU[1](H, a))
  \]
  and therefore bijective. Showing the universal property means
  showing that $\inv\Yoneda$ is multi-natural and that:
  \begin{center}
    $\quiver{axiom 3b proof goal}$
  \end{center}
  Multi-naturality follows from the Isomorphism Lemma, and we'll
  establish the equation by direct calculation, taking any $z \in
  \Obj{\UU[3]}, \zvec \in \opposite{\UU[3]}$ and $w : (\zvec) \to z$
  in $\UU[3]$ and showing that:
  \begin{multline}\label{axiom3b RTP}
  \parent{(\Psh G \hcompose (\map H) \hcompose G)\vcompose (\map (\yoneda \compose G))}_{z, \seq[k=1][1]\zvec} w
  \\=
  \parent{
    (\inv\Yoneda \hcompose (G \compose H)) \vcompose
    (\map(H\compose G))}_{z, \seq[k=1][1]\zvec} w
  \mspace{100mu}
  \end{multline}
  Type-checking, this equation is between natural multi-transformations within the multi-hom-set:
  \[
  \Psh{\UU[2]}(\yoneda (G[\zvec]), \UU(H, HGz))
  \]
  and so let $n \definedby \abs\zvec$, and take an $n$-ary partition
  $\partit$ of some $\yvec \in \opposite{\UU[2]}$, and any $n$-tuple
  of multi-morphisms:
  \[
  \seq[i=1][n]{(\partit_i) \xto{v_i} G\zvec_i)} \in \prod_{i=1}^n\yoneda(G\zvec_i)(\partit_i)
  \]
  and apply both sides of \eqref{axiom3b RTP} at the $\partit$-component to this tuple.

  On the LHS:
  \begin{multline*}
    \parent{\parent{(\Psh G \hcompose (\map H) \hcompose G)\vcompose (\map (\yoneda \compose G))}_{z, \seq[k=1][1]\zvec} w}_{\partit}\seq[i=1][n]{v_i}
    \\
    \begin{aligned}
      &{}=
      \parent{(\Psh G (\map H)_{Gz})_{\seq[k=1][1]\zvec} (\yoneda (Gw))}_{\partit}\seq[i=1][n]{v_i}
      \explain={$\Psh G$ action on natural \\multi-transformations}
      (\map H \compose (\yoneda (Gw)))_{\partit}\seq[i=1][n]{v_i}
      \\&{}=
      \map H(Gw \compose \seq[i=1][n]{v_i})
      \explain={$\map H$-def, \\multi-functoriality}
      HGw \compose \seq[i=1][n]{Hv_i}
    \end{aligned}
  \end{multline*}
  On the RHS:
  \begin{multline*}
    \parent{\parent{
    (\inv\Yoneda \hcompose (G \compose H)) \vcompose
        (\map(H\compose G))}_{z, \seq[k=1][1]\zvec} w}_{\partit}\seq[i=1][n]{v_i}
    \\
    \begin{aligned}
      &{}=
      \inv\Yoneda_{;\UU(H;a);\yvec}\parent{\map (H \compose G)w}\seq[i=1][n]{v_i}
      \explain={$\inv\Yoneda$-def,$\map$-def}
      \UU(H[\seq[i=1][n]{v_i}];a)(HGw)
      =
      HGw \compose \seq[i=1][n]{Hv_i}
    \end{aligned}
  \end{multline*}
  and so \eqref{axiom3b RTP} holds, and $\theta_{G,H}$ is a
  multi-natural isomorphism.
\end{proof}

% \subsection{Failure of Axiom $3^*$}
% Street and Walters also suggest a stronger axiom:
% \paragraph{Axiom $3^*$}
% Every multi-natural transformation $\sigma$ as on the left such that
% the diagram on the right is an absolute lifting of $\yoneda$ through
% $H$, must be a multi-natural isomorphism:
% \begin{center}
%   $\quiver{axiom 3* assumption}$
% \end{center}
% \begin{counterexample}\emph{Failure of Axiom $3^*$.}
%   Take:
%   \begin{itemize}
%   \item
%     $\UU$ to be the multi-category with no objects and no
%     multi-morphisms whatsoever;
%   \item $\UU[2]$ to be the multi-category with a single object $x$ and
%     only the identity on $x$;
%   \item $F : \UU \to \UU[2]$ to be the (unique) multi-functor of this type.
%   \end{itemize}
%   The relative Yoneda functor $\UU[2](F,-)$ then sends $x$ to the
%   presheaf constantly taking the empty set.
%   \begin{itemize}
%   \item Take $H : \UU[2] \to \Psh\UU$ to be the presheaf constantly
%     taking the singleton $\terminal \definedby \set{\star}$.
%   \item Take $\sigma : \UU[2](F, -) \to H$ to be the only possible
%     multi-transformation.
%   \end{itemize}
%   Since the sink in $\sigma$'s naturality diagrams is the singleton,
%   it is a natural multi-transformation $\sigma : \UU[2](F, -) \to
%   H$. By the Isomorphism Lemma, since the source and target sets of
%   $\sigma_{x, \seq{}} : \emptyset \to \terminal$ have different
%   cardinalities, $\sigma$ is not a natural multi-isomorphism.

%   However, the appropriate triangle is an absolute lifting. Indeed,
%   take any multifunctors: $W : \UU[3] \to \UU$ and $G : \UU[3] \to
%   \UU[2]$. We will show that there are unique multi-natural
%   transformations $\alpha$ (left) and $\beta$ (right):
%   \begin{center}
%     $\quiver{Axiom 3* counterexample goal}$
%   \end{center}
%   The pasting (on the right) then forms the bijection we want.

%   Since $\UU$ has no objects or morphisms, we must have that $\UU[3]$
%   doesn't have any objects or morphisms either, and so it is identical
%   $\UU$. Since $F$ is the unique multi-functor out of $\UU$, we have
%   that $G$ is the unique multi-functor out of $\UU[3]$, and $\beta$
%   must be the identification of the two. Since $H \compose G$ is
%   componentwise a singleton, each component of $\alpha$ must be the
%   unique map into the singleton, and so $\alpha$ is unique.
% \end{counterexample}

\subsection{Representability}
We can now reap the fruit of our labour:

\begin{proposition}\emph{Characterisation: multi-categorical adjunctions.}
\propositionlabel{representability multi-categorical adjunctions}
Let $\UU[1]$ be a locally-small multi-category. TFAE:
\begin{itemize}
\item An adjunction $\triple{\quiver{adjunction type}}{\Id \xto{\unit}
  UF}{FU \xto\counit \Id}$ in $\MultiCAT$ with $F$ admissible.
\item A multi-natural isomorphism $\quiver{adjunction hom-set bijection}$ with $F$ admissible.
\item A multi-functor $U : \UU[2] \to \UU[1]$, a choice of objects $Fa
  \in \UU[2]$ for all $a \in \UU[1]$ together with a (multi-)morphism
  $\unit_a : (a) \to UFa$ such that:
  \begin{itemize}
  \item All homs $\UU[2](Fa_1, \ldots, Fa_n; x)$ are small; and
  \item The components of the transformation $\unit$ are
    multi-universal: for all $(a_1, \ldots, a_n) \xto f U x$ in
    $\UU[1]$, there exist a unique $(Fa_1, \ldots, Fa_n) \xto{\psi(f)}
    x$ in $\UU[2]$ such that:
    \[
    U (\psi(f))\compose (\unit_{a_1}, \ldots, \unit_{a_n}) = f
    \]
  \end{itemize}
\item
  An admissible multi-functor $F : \UU[1] \to \UU[2]$, a choice of
  objects $Ux \in \UU[1]$ for all $x \in \UU[2]$ together with a
  (multi-)morphism $\counit_x : (FUx) \to x$ that is universal: for
  all $(Fa_1, \ldots, Fa_n) \xto g x$ in $\UU[2]$, there exist a
  unique $(a_1, \ldots, a_n) \xto{\phi g} U x$ in $\UU[1]$ such that:
  \[
  \counit_x\compose (F (\phi g)) = g
  \]
\end{itemize}
\end{proposition}
\begin{proof}
  The first two items are Prop.~8 of \cite{street-walters-yoneda-structures}.
  Given a multi-natural isomorphism $\phi$ with inverse $\psi$, take:
  \[
  \unit_a \definedby \phi\id
  \qquad
  \counit_x \definedby \psi\id
  \]
  Take any $f : (\avec) \to U x$, $n \definedby \abs\avec$, and chase
  $\seq[i=1][n]{\id[a_n]} \in \prod_{i=1}^n\UU[2](Fa_i; Fa_i)$ around
  the multi-naturality condition for $\phi$ when instantiated for
  $\psi f$ in $\UU[2]$.  OTOH:
  \[
  \phi\parent{\UU[2](F;\psi f)\seq[i=1][n]{\id[a_i]}}
  =
  \phi\parent{\psi f \compose \seq[i=1][n]{\id[a_i]}}
  =
  \phi\parent{\psi f}
  =
  f
  \]
  and OTOH:
  \[
  \UU[1](\Id; U(\psi f))\seq[i=1][n]{\phi(\id[a_i])}
  =
  U(\psi f)\compose\seq[i=1][n]{\unit_{a_i}}
  \]
  and we proved the third item.

  Take any $g : (F[\avec]) \to x$ in $\UU[2]$, and chase $\id[x]$
  around the naturality condition for $\psi_x$ when applied to $\phi g
  : \seq[k=1][1]{U x} \from \avec$. OTOH:
  \[
  \psi\parent{\UU(\phi g; Ux)\id}
  =
  \psi\parent{\id \compose (\phi g)}
  = \psi (\phi g) = g
  \]
  and OTOH:
  \[
  \UU[2](F(\phi g); U x)(\psi\id[x])
  =
  \counit_x\compose (F\phi g)
  \]
  as we wanted.

  Therefore, the first two conditions imply the last two conditions.
  We'll now show that the third and fourth conditions imply the third.
  The exists-unique conditions construct the required bijections, and
  we need to equip $F$ or $U$ with a multi-functorial action and show
  that the given bijection is multi-natural in $x$ and natural in
  $\avec$, in each case.

  First, assume we have a choice of multi-universal units. Extend $F$
  to a functor, by sending each $f : (\avec) \to b$ in $\UU[1]$ to:
  \[
  \psi(\unit_b\compose (f))
  \qquad
  \text{i.e., the unique $g : (F[\avec]) \to Fb$ satisfying: }
  \qquad
  Ug \compose \seq[i=1][\abs\avec]{\unit_{\avec_i}} = \unit_b \compose (f)
  \]
  This assignment is functorial:
  \[
  U\id[Fa] \compose (\unit_{a})
  \explain={$U$ multi-functor}
  \unit_{a}
  =
  \unit_{a}\compose(\id[a])
  \qquad\implies\qquad
  F\id[a] = \id[Fa]
  \]
  and for all composable $f : (\bvec) \to c$, $n \definedby \abs\bvec$
  and $\fvec \definedby \seq[i=1][n]{f_i : (\avec_{\pair i*}) \to
    \bvec_i}$, $p \definedby \abs\avec$,$m_i \definedby
  \abs{\avec_{\pair i *}}$, we have:
  \begin{align*}
    U(Ff\compose \seq[i=1][n]{F\fvec_i}) \compose \seq[\ell = 1][p]{\unit_{a_\ell}}
    &{}\explain={$U$ multi-functor}
    (UFf \compose UF\fvec_i)\compose \seq[\ell = 1][p]{\unit_{a_\ell}}
    \explain={$\compose$-assoc.}
    UFf \compose (UF\fvec_i\compose \seq[j = 1][m_i]{\unit_{a_{\pair ij}}})
    \\&{}= UFf \compose (\unit_{\bvec_i}\compose (\fvec_i))
    = (UFf \compose \unit_{\bvec_i})\compose (\fvec)
    = (\unit_c \compose UFf) \compose (\fvec)
    \\&{}= \unit_c \compose (UFf \compose (\fvec))
  \end{align*}
  and so $F(f \compose (\fvec)) = F f \compose (F[\fvec])$.

  Fix $x$. To show that $\psi_x$ is natural, take any $\avec \xfrom \uvec
  \bvec$ in $\UU[1]$, and chasing any $f : (\avec) \to U x$ around the
  naturality condition for $\psi$ gives us the proof obligation:
  \[
  \psi(f \compose (\uvec))
  =
  \psi f \compose (F[\uvec])
  \]
  Calculate:
  \begin{align*}
    U\parent{\psi f \compose (F[\uvec])} \compose
    \seq[\ell = 1][\abs\bvec]{\unit_{\avec_\ell}}
    &{}\explain={$F$-multi-functor\\$\compose$-assoc.}
    U(\psi f) \compose \seq[i=1][\abs\avec]{UF\uvec_i \compose
      \seq[j = 1][\abs{\bvec_{\pair i*}}]{\unit_{\avec_{\pair ij}}}}
    \explain={$F$-def}
    U(\psi f) \compose \seq[i=1][\abs\avec]{\unit_{\bvec}\compose (\uvec_i)}
    \\&\explain={$\compose$-assoc\\$\psi$-def}
    f \compose (\uvec)
  \end{align*}
  and so $\psi_x$ is natural.

  To show $\phi$ is multi-natural, take any $(\xvec)\xto v y$ in
  $\UU[2]$, and chase an arbitrary tuple
  $\seq[i=1][\abs\xvec]{\avec_{\pair i*}\xto{\fvec_i} U\xvec_i}$ in
  the multi-naturality condition to get the proof obligation:
  \[
  \psi(Uv \compose (\fvec)) = v \compose (\psi[\fvec])
  \qquad
  \impliedby
  \qquad
      U\parent{v \compose (\psi[\fvec])} \compose
    \seq[\ell=1][\abs\avec]{\unit_{\avec_\ell}}
    \explain={$U$ multi functor\\$\compose$-assoc.\\$\psi$-def}
    Uv \compose (\seq[i=1][\abs\xvec]{\fvec_i}
  \]
  and so $\psi$ is multi-natural. By the Isomorphism Lemma, it is a
  multi-natural isomorphism, and we proved that the third condition
  implies the second.

  The proof that the final condition implies the second is similar.
  To extend $U$ to a functor, send $g : (\xvec)\to y$ in $\UU[2]$ to:
  \[
  \phi(g\compose\seq[i=1][\abs\xvec]{\counit_{\xvec_i}})
  \qquad
  \text{i.e., the unique $f : (U[\xvec]) \to Uy$ satisfying: }
  \qquad
  \counit_y\compose (Ff) = f\compose\seq[i=1][\abs\xvec]{\counit_{\xvec_i}}
  \]
  This assignment is functorial:
  \[
  \counit \compose (F\id) \explain={$F$ multi-functor}
  \counit \compose (\id) = \id\compose (\counit)
  \qquad\implies\qquad
  U\id = \id
  \]
  and for all composable $(\xvec) \xto{\gvec} (\yvec) \xto g z$:
  \[
  \counit_z \compose(F(Ug \compose (U[\gvec])))
  \explain={$F$ multi-functor\\$\compose$ assoc.}
  (\counit_z\compose FUg)\compose (FU[\gvec])
  =
  g \compose \seq[i]{\counit_{\yvec_i}\compose (FU[\gvec_i])}
  =
  (g \compose (\gvec))\compose \seq[\ell]{\counit_{\xvec_\ell}}
  \]
  and so $U(g \compose (\gvec)) = Ug \compose (U[\gvec])$.

  Fix $x$. To show that $\phi_x$ is natural, take any $\avec
  \xfrom\uvec \bvec$ in $\UU[1]$ and chase any $v : (F[\avec]) \to x$
  through the multi-naturality condition to get the proof obligation:
  \[
  \phi\parent{v \compose F[\uvec]} =
  (\phi v) \compose (\uvec)
  \impliedby
  \counit_x\compose F((\phi v)\compose (\uvec))
  \explain={$F$ multi-functor\\$\compose$ assoc.}
  (\counit_x\compose F\phi v))\compose (F[\uvec])
  \explain={$\phi$-def}
  v \compose (F[\uvec])
  \]

  Finally, to show that $\phi$ is natural, take any $v : (\xvec) \to
  y$ and chase $\fvec : \avec \to U[\xvec]$ in the multi-naturality
  condition to get the proof obligation:
  \[
  \phi\parent{v \compose (\fvec)} = Uv \compose \phi[\fvec]
  \impliedby
  \begin{aligned}[t]
  \counit_y\compose (F\parent{Uv \compose \phi[\fvec]})
  &{}\explain={$F$ multi functor\\$\compose$ assoc.}
  (\counit_y \compose FUv)\compose (F\phi[\fvec])
  \explain={$U$-def}
  v \compose \seq[i]{\counit_{\xvec_i}\compose F\phi[\fvec]}
  \\&{}\explain={$\phi$-def}
  v \compose (\fvec)
  \end{aligned}
  \]
  as we wanted.
\end{proof}
